\documentclass[10pt,a4paper]{amsart}
\usepackage[margin=19mm]{geometry}
\usepackage{amsmath,amssymb,mathtools}
\usepackage[scr=rsfs,cal=euler]{mathalfa}
\usepackage{xfrac}
\usepackage[unicode,hidelinks,bookmarks=false]{hyperref}
\newcommand{\ie}{i.e.\ }
\newcommand{\cf}{cf.\ }
\newcommand{\resp}{respectively\ }
\newcommand{\Subsection}{\subsection}
\providecommand{\nobreakdash}{\nobreak\hbox{-}}
\setlength{\emergencystretch}{2em}
\multlinegap0pt
\usepackage{bm}
\usepackage{bbm}
\usepackage{dsfont}
\usepackage{xspace}
\usepackage{cancel}
\usepackage{mathtools}
\usepackage{amsthm}
\makeatletter
\@ifundefined{theorem}{\newtheorem{theorem}{Theorem}}{}
\@ifundefined{lemma}{\newtheorem{lemma}[theorem]{Lemma}}{}
\@ifundefined{remark}{\newtheorem{remark}[theorem]{Remark}}{}
\makeatother
\title[Integral representations and asymptotic behaviors of the Mul]{Integral representations and asymptotic behaviors of the Multivariate Mittag-Leffler function}
\author{Damir Shamuratov}
\date{Source: arXiv:2607.19103v1 (CC BY 4.0)}
\begin{document}
\maketitle

\noindent\emph{Source and licence.} Integral representations and asymptotic behaviors of the Multivariate Mittag-Leffler function. Version arXiv:2607.19103v1 (CC BY 4.0), distributed under the Creative Commons Attribution 4.0 International licence, as recorded in the arXiv licence metadata for this exact version and shown on the abstract page https://arxiv.org/abs/2607.19103v1. Full rights notice: licenses/arxiv-2607.19103-CC-BY-4.0.txt.\par\smallskip

\noindent\textit{Editorial scope.} Counted unit: the Lemma whose source label is \texttt{l3.2} in the article (source0.tex lines 614--658, source label \texttt{l3.2}), delivered together with the article's own complete original proof of it (source0.tex lines 659--1114, one \texttt{proof} environment of 456 lines). No mathematics is written, repaired or re-derived: statement and proof are the article's own text line for line. Required context retained uncounted: 3.1 (lines 426--432); 1.3 (lines 549--564). Every definition, display and label of those lines is retained unchanged. Omitted from this package and not redistributed: 1--425, 433--548, 565--613, 1115--2408. Nothing inside a retained excerpt is omitted or altered: every retained body is delivered exactly as the article's lines are written, so no omission receipt of any kind accompanies this package. Citations: the retained excerpts carry no citation command at all, so no bibliography entry of the article is transcribed and no reference is redistributed. Reference adaptation: none was needed and none was made; every reference inside the delivered unit resolves inside it, so no unresolved reference or citation is produced. Printed slips kept literal: no printed word, symbol, hypothesis or number of the counted statement or of its proof was altered. Layout and class: the article's own class is \texttt{amsart} with packages including amsmath, amsthm, amscd, amsfonts, amssymb, graphicx, color, mathrsfs, hyperref, tikz; the wrapper is the lane's pinned \texttt{amsart} editorial wrapper at 10pt on A4, which restates the theorem-like environments the retained text needs and adds the article's own macro definitions that the retained text needs (none), with extra wrapper packages (bm, bbm, dsfont, xspace, cancel, mathtools). The delivered numbering is the unit's own. Assets: no retained span contains a figure environment, a graphics-inclusion command, a PDF-page inclusion, a table or a TikZ picture; no image file is copied into this package, the package declares no asset dependency and no image file is redistributed with it. Licence: version arXiv:2607.19103v1 of this article is distributed under the Creative Commons Attribution 4.0 International licence, as recorded in the arXiv licence metadata for this exact version and shown on the abstract page https://arxiv.org/abs/2607.19103v1. A full rights notice is delivered at licenses/arxiv-2607.19103-CC-BY-4.0.txt. Characters: every retained excerpt is pure ASCII, so the delivered engine, XeLaTeX with its default Latin Modern text font, reports no missing character.
\begin{equation}\label{3.1}
\frac{\pi\alpha\beta\gamma}{2}
<
\eta
\leq
\min\!\bigl(\pi,\pi\alpha\beta\gamma\bigr).
\end{equation}

\begin{equation}\label{1.3}
\begin{aligned}
E_{\alpha,\beta,\gamma}(x,y,z;\mu)
&=
\frac{1}{2\pi i}\,
\frac{1}{\gamma}
\int_{\omega(\varepsilon_\gamma;\eta_\gamma)}
\frac{
e^{s^{1/\gamma}}
\,s^{\frac{\alpha+\beta+1-\mu}{\gamma}}
}
{
(s^{\alpha/\gamma}-x)(s^{\beta/\gamma}-y)(s-z)
}
\,ds.
\end{aligned}\end{equation}

\begin{lemma}\label{l3.2}
     Let $0 < \alpha, \beta,\gamma < 2$ and $\alpha\beta\gamma < 2$. Let $\mu$ be any complex number and let $\eta$ satisfy the condition \eqref{3.1}.
If
$
x \in \Omega^{(+)}({\varepsilon}_\alpha;\eta_\alpha)
$,\,\, 
$y \in \Omega^{(+)}({\varepsilon}_\beta;\eta_\beta),
$ and $z \in \Omega^{(+)}({\varepsilon}_\gamma;\eta_\gamma)
$,
where
$
{\varepsilon}_\alpha:={\varepsilon}^{\frac{1}{\beta\gamma}},\,
{\varepsilon}_\beta:={\varepsilon}^{\frac{1}{\alpha\gamma}}, \,{\varepsilon}_\gamma:={\varepsilon}^{\frac{1}{\alpha\beta}},\,
\eta_\alpha:=\cfrac{{\eta}}{\beta\gamma},\,
\eta_\beta:=\cfrac{{\eta}}{\alpha\gamma},\,
\eta_\gamma:=\cfrac{{\eta}}{\alpha\beta},
$
then the following Hankel integral representation holds:
\begin{equation}\label{3.9}\begin{aligned}
&E_{\alpha,\beta,\gamma}(x,y,z;\mu)\\
&=\frac{1}{\alpha}\frac{
e^{x^{1/\alpha}}
\,x^{\frac{\beta+\gamma+1-\mu}{\alpha}}}
{(x^{\beta/\alpha}-y)(x^{\gamma/\alpha}-z)}+\frac{1}{\beta}\frac{
e^{y^{1/\beta}}
\,y^{\frac{\alpha+\gamma+1-\mu}{\beta}}}
{(y^{\alpha/\beta}-x)(y^{\gamma/\beta}-z)}+\frac{1}{\gamma}
\frac{
e^{z^{1/\gamma}}
\,z^{\frac{\alpha+\beta+1-\mu}{\gamma}}}
{(z^{\alpha/\gamma}-x)(z^{\beta/\gamma}-y)}\\
&+
\frac{1}{2\pi i}\,
\frac{1}{\alpha\beta\gamma}
\int_{\omega(\varepsilon;\eta)}
\frac{
e^{s^{1/(\alpha\beta\gamma)}}
\,s^{\frac{\alpha+\beta+\gamma+1-\mu}{\alpha\beta\gamma}-1}
}
{
(s^{1/\beta\gamma}-x)(s^{1/\alpha\gamma}-y)(s^{1/\alpha\beta}-z)
}
\,ds.\end{aligned}
\end{equation}
\end{lemma}

\begin{proof}
   First, suppose that only the point $z$ lies to the right of the Hankel contour
$\omega(\varepsilon_\gamma;\eta_\gamma)$, namely,
$
z\in\Omega^{(+)}(\varepsilon_\gamma;\eta_\gamma).
$
Choose an arbitrary number $\widetilde{\varepsilon}_{\gamma}$ satisfying
$\widetilde{\varepsilon}_{\gamma}>|z|$. Then
$
z\in\Omega^{(-)}(\widetilde{\varepsilon}_{\gamma};\eta_\gamma).
$
Furthermore, setting
$
\widetilde{\varepsilon}_{\alpha}
=\widetilde{\varepsilon}_{\gamma}^{\frac{\alpha}{\gamma}},
\,\,
\widetilde{\varepsilon}_{\beta}
=\widetilde{\varepsilon}_{\gamma}^{\frac{\beta}{\gamma}},
$
it follows that
$
x\in\Omega^{(-)}(\widetilde{\varepsilon}_{\alpha};\eta_\alpha),
\,\,
y\in\Omega^{(-)}(\widetilde{\varepsilon}_{\beta};\eta_\beta).
$
Hence, applying the integral representation \eqref{1.3}, we arrive at
\begin{equation}\label{3.10}E_{\alpha,\beta,\gamma}(x,y,z;\mu)
=
\frac{1}{2\pi i}\,
\frac{1}{\gamma}
\int_{\omega(\widetilde{\varepsilon}_{\gamma};\eta_\gamma)}
\frac{
e^{s^{1/\gamma}}
\,s^{\frac{\alpha+\beta+1-\mu}{\gamma}}
}
{
(s^{\alpha/\gamma}-x)(s^{\beta/\gamma}-y)(s-z)
}
\,ds.\end{equation}

Moreover, if
$
\varepsilon_\gamma<|z|<\widetilde{\varepsilon}_{\gamma},
$
then
$
|\arg z|<\eta_\gamma,
$
and, by Cauchy's theorem,
\begin{equation}\label{3.11}\begin{aligned}
E_{\alpha,\beta,\gamma}(x,y,z;\mu)
&=
\frac{1}{2\pi i}\,
\frac{1}{\gamma}
\int_{\omega(\widetilde{\varepsilon}_{\gamma};\eta_\gamma)
-
\omega(\varepsilon_\gamma;\eta_\gamma)}
\frac{
e^{s^{1/\gamma}}
\,s^{\frac{\alpha+\beta+1-\mu}{\gamma}}
}
{
(s^{\alpha/\gamma}-x)(s^{\beta/\gamma}-y)(s-z)
}ds\\
&=\frac{1}{\gamma}
\frac{
e^{z^{1/\gamma}}
\,z^{\frac{\alpha+\beta+1-\mu}{\gamma}}}
{(z^{\alpha/\gamma}-x)(z^{\beta/\gamma}-y)}\end{aligned}
\end{equation}

Therefore, from \eqref{3.10} and \eqref{3.11}, we obtain the following integral representation 
\begin{equation}\label{3.12}\begin{aligned}
E_{\alpha,\beta,\gamma}(x,y,z;\mu)
&=\frac{1}{\gamma}\frac{
e^{z^{1/\gamma}}
\,z^{\frac{\alpha+\beta+1-\mu}{\gamma}}}
{(z^{\alpha/\gamma}-x)(z^{\beta/\gamma}-y)}\\
&+
\frac{1}{2\pi i}\,
\frac{1}{\alpha\beta\gamma}
\int_{\omega(\varepsilon;\eta)}
\frac{
e^{s^{1/(\alpha\beta\gamma)}}
\,s^{\frac{\alpha+\beta+\gamma+1-\mu}{\alpha\beta\gamma}-1}
}
{
(s^{1/\beta\gamma}-x)(s^{1/\alpha\gamma}-y)(s^{1/\alpha\beta}-z)
}
\,ds.\end{aligned}\end{equation}

\begin{remark} Analogously, for  $x\in\Omega^{(+)}(\varepsilon_{\alpha};\eta_\alpha),$
\,
$y\in\Omega^{(-)}(\varepsilon_{\beta};\eta_\beta)$,\, $z\in\Omega^{(-)}(\varepsilon_{\gamma};\eta_\gamma)$ or $x\in\Omega^{(-)}(\varepsilon_{\alpha};\eta_\alpha),$
\,
$y\in\Omega^{(+)}(\varepsilon_{\beta};\eta_\beta)$,\, $z\in\Omega^{(-)}(\varepsilon_{\gamma};\eta_\gamma)$
we obtain the following integral representations, respectively 
$$\begin{aligned}
E_{\alpha,\beta,\gamma}(x,y,z;\mu)
&=\frac{1}{\alpha}\frac{
e^{x^{1/\alpha}}
\,x^{\frac{\beta+\gamma+1-\mu}{\alpha}}}
{(x^{\beta/\alpha}-y)(x^{\gamma/\alpha}-z)}\\
&+
\frac{1}{2\pi i}\,
\frac{1}{\alpha\beta\gamma}
\int_{\omega(\varepsilon;\eta)}
\frac{
e^{s^{1/(\alpha\beta\gamma)}}
\,s^{\frac{\alpha+\beta+\gamma+1-\mu}{\alpha\beta\gamma}-1}
}
{
(s^{1/\beta\gamma}-x)(s^{1/\alpha\gamma}-y)(s^{1/\alpha\beta}-z)
}
\,ds,\end{aligned}
$$
and 
$$\begin{aligned}
E_{\alpha,\beta,\gamma}(x,y,z;\mu)
&=\frac{1}{\beta}\frac{
e^{y^{1/\beta}}
\,y^{\frac{\alpha+\gamma+1-\mu}{\beta}}}
{(y^{\alpha/\beta}-x)(y^{\gamma/\beta}-z)}\\
&+
\frac{1}{2\pi i}\,
\frac{1}{\alpha\beta\gamma}
\int_{\omega(\varepsilon;\eta)}
\frac{
e^{s^{1/(\alpha\beta\gamma)}}
\,s^{\frac{\alpha+\beta+\gamma+1-\mu}{\alpha\beta\gamma}-1}
}
{
(s^{1/\beta\gamma}-x)(s^{1/\alpha\gamma}-y)(s^{1/\alpha\beta}-z)
}
\,ds.\end{aligned}
$$\end{remark}

Second, suppose that both $y$ and $z$ lie to the right of the corresponding Hankel contours, that is,
$
y\in\Omega^{(+)}(\varepsilon_\beta;\eta_\beta),
\,\,
z\in\Omega^{(+)}(\varepsilon_\gamma;\eta_\gamma).
$
From \eqref{3.12}, it remains to analyze the remaining contour integral. Applying the change of variables
$
s^{1/\alpha\gamma}=t,
$
we obtain
\begin{equation}\label{3.13}\begin{aligned}
E_{\alpha,\beta,\gamma}(x,y,z;\mu)
&=\frac{1}{\gamma}\frac{
e^{z^{1/\gamma}}
\,z^{\frac{\alpha+\beta+1-\mu}{\gamma}}}
{(z^{\alpha/\gamma}-x)(z^{\beta/\gamma}-y)}\\
&+
\frac{1}{2\pi i}\,
\frac{1}{\beta}
\int_{\omega(\varepsilon_\beta;\eta_\beta)}
\frac{
e^{t^{1/\beta}}
\,t^{\frac{\alpha+\gamma+1-\mu}{\beta}}
}
{
(t^{\alpha/\beta}-x)(t-y)(t^{\gamma/\beta}-z)
}
\,dt.\end{aligned}\end{equation}
Under the transformation
$
s=t^{\alpha\gamma},
$
the Hankel contour
$
\omega(\varepsilon;\eta)
$
is mapped onto the Hankel contour
$
\omega(\varepsilon_\beta;\eta_\beta).
$ Choose an arbitrary number $\widetilde{\varepsilon}_{\beta}$ satisfying
$\widetilde{\varepsilon}_{\beta}>|y|$. Then
$
y\in\Omega^{(-)}(\widetilde{\varepsilon}_{\beta};\eta_\beta).
$
Furthermore, setting
$
\widetilde{\varepsilon}_{\alpha}
=\widetilde{\varepsilon}_{\beta}^{\frac{\alpha}{\beta}},
\,\,
\widetilde{\varepsilon}_{\gamma}
=\widetilde{\varepsilon}_{\beta}^{\frac{\gamma}{\beta}},
$
it follows that
$
x\in\Omega^{(-)}(\widetilde{\varepsilon}_{\alpha};\eta_\alpha),
\,\,
z\in\Omega^{(-)}(\widetilde{\varepsilon}_{\gamma};\eta_\gamma).
$ Hence, applying the integral representation \eqref{3.13}, we arrive at
\begin{equation}\label{3.14}\begin{aligned}
E_{\alpha,\beta,\gamma}(x,y,z;\mu)
&=\frac{1}{\gamma}\frac{
e^{z^{1/\gamma}}
\,z^{\frac{\alpha+\beta+1-\mu}{\gamma}}}
{(z^{\alpha/\gamma}-x)(z^{\beta/\gamma}-y)}\\
&+
\frac{1}{2\pi i}\,
\frac{1}{\beta}
\int_{\omega(\widetilde{\varepsilon}_{\beta};\eta_\beta)}
\frac{
e^{t^{1/\beta}}
\,t^{\frac{\alpha+\gamma+1-\mu}{\beta}}
}
{
(t^{\alpha/\beta}-x)(t-y)(t^{\gamma/\beta}-z)
}
\,dt.\end{aligned}\end{equation}
Moreover, if
$
\varepsilon_\beta<|y|<\widetilde{\varepsilon}_{\beta},
$
then
$
|\arg y|<\eta_\beta,
$
and, by Cauchy's theorem,
\begin{equation}\label{3.15}\begin{aligned}
E_{\alpha,\beta,\gamma}(x,y,z;\mu)
&=
\frac{1}{2\pi i}\,
\frac{1}{\beta}
\int_{\omega(\widetilde{\varepsilon}_{\beta};\eta_\beta)-\omega(\varepsilon_{\beta};\eta_\beta)}
\frac{
e^{t^{1/\beta}}
\,t^{\frac{\alpha+\gamma+1-\mu}{\beta}}
}
{
(t^{\alpha/\beta}-x)(t-y)(t^{\gamma/\beta}-z)
}
\,dt\\
&=\frac{1}{\beta}\frac{
e^{y^{1/\beta}}
\,y^{\frac{\alpha+\gamma+1-\mu}{\beta}}}
{(y^{\alpha/\beta}-x)(y^{\gamma/\beta}-z)}\end{aligned}
\end{equation}
Combining  \eqref{3.15}, \eqref{3.14}, \eqref{3.13} and \eqref{3.12}, we obtain the following integral representation 
\begin{equation}\label{3.16}\begin{aligned}
E_{\alpha,\beta,\gamma}(x,y,z;\mu)
&=\frac{1}{\gamma}\frac{
e^{z^{1/\gamma}}
\,z^{\frac{\alpha+\beta+1-\mu}{\gamma}}}
{(z^{\alpha/\gamma}-x)(z^{\beta/\gamma}-y)}+\frac{1}{\beta}\frac{
e^{y^{1/\beta}}
\,y^{\frac{\alpha+\gamma+1-\mu}{\beta}}}
{(y^{\alpha/\beta}-x)(y^{\gamma/\beta}-z)}\\
&+
\frac{1}{2\pi i}\,
\frac{1}{\alpha\beta\gamma}
\int_{\omega(\varepsilon;\eta)}
\frac{
e^{s^{1/(\alpha\beta\gamma)}}
\,s^{\frac{\alpha+\beta+\gamma+1-\mu}{\alpha\beta\gamma}-1}
}
{
(s^{1/\beta\gamma}-x)(s^{1/\alpha\gamma}-y)(s^{1/\alpha\beta}-z)
}
\,ds.\end{aligned}\end{equation}

\begin{remark} Analogously, for  $x\in\Omega^{(+)}(\varepsilon_{\alpha};\eta_\alpha),$
\,
$y\in\Omega^{(+)}(\varepsilon_{\beta};\eta_\beta)$,\, $z\in\Omega^{(-)}(\varepsilon_{\gamma};\eta_\gamma)$ or $x\in\Omega^{(+)}(\varepsilon_{\alpha};\eta_\alpha),$
\,
$y\in\Omega^{(-)}(\varepsilon_{\beta};\eta_\beta)$,\, $z\in\Omega^{(+)}(\varepsilon_{\gamma};\eta_\gamma)$
we obtain the following integral representations, respectively 
$$\begin{aligned}
E_{\alpha,\beta,\gamma}(x,y,z;\mu)
&=\frac{1}{\alpha}\frac{
e^{x^{1/\alpha}}
\,x^{\frac{\beta+\gamma+1-\mu}{\alpha}}}
{(x^{\beta/\alpha}-y)(x^{\gamma/\alpha}-z)}+\frac{1}{\beta}\frac{
e^{y^{1/\beta}}
\,y^{\frac{\alpha+\gamma+1-\mu}{\beta}}}
{(y^{\alpha/\beta}-x)(y^{\gamma/\beta}-z)}\\
&+
\frac{1}{2\pi i}\,
\frac{1}{\alpha\beta\gamma}
\int_{\omega(\varepsilon;\eta)}
\frac{
e^{s^{1/(\alpha\beta\gamma)}}
\,s^{\frac{\alpha+\beta+\gamma+1-\mu}{\alpha\beta\gamma}-1}
}
{
(s^{1/\beta\gamma}-x)(s^{1/\alpha\gamma}-y)(s^{1/\alpha\beta}-z)
}
\,ds,\end{aligned}
$$
and 
$$\begin{aligned}
E_{\alpha,\beta,\gamma}(x,y,z;\mu)
&=\frac{1}{\alpha}\frac{
e^{x^{1/\alpha}}
\,x^{\frac{\beta+\gamma+1-\mu}{\alpha}}}
{(x^{\beta/\alpha}-y)(x^{\gamma/\alpha}-z)}+\frac{1}{\gamma}\frac{
e^{z^{1/\gamma}}
\,z^{\frac{\alpha+\beta+1-\mu}{\gamma}}}
{(z^{\alpha/\gamma}-x)(z^{\beta/\gamma}-y)}\\
&+
\frac{1}{2\pi i}\,
\frac{1}{\alpha\beta\gamma}
\int_{\omega(\varepsilon;\eta)}
\frac{
e^{s^{1/(\alpha\beta\gamma)}}
\,s^{\frac{\alpha+\beta+\gamma+1-\mu}{\alpha\beta\gamma}-1}
}
{
(s^{1/\beta\gamma}-x)(s^{1/\alpha\gamma}-y)(s^{1/\alpha\beta}-z)
}
\,ds.\end{aligned}
$$\end{remark}

Third, all $x$, $y$ and $z$ points lie to the right of the corresponding Hankel contours, that is,
$
y\in\Omega^{(+)}(\varepsilon_\beta;\eta_\beta),
\,\,
z\in\Omega^{(+)}(\varepsilon_\gamma;\eta_\gamma),\,\,x\in\Omega^{(+)}(\varepsilon_\alpha;\eta_\alpha)
$.
From \eqref{3.16}, it remains to analyze the remaining contour integral. Applying the change of variables
$
s^{1/\beta\gamma}=t,
$
we obtain
\begin{equation}\label{3.17}\begin{aligned}
E_{\alpha,\beta,\gamma}(x,y,z;\mu)
&=\frac{1}{\gamma}\frac{
e^{z^{1/\gamma}}
\,z^{\frac{\alpha+\beta+1-\mu}{\gamma}}}
{(z^{\alpha/\gamma}-x)(z^{\beta/\gamma}-y)}+\frac{1}{\beta}\frac{
e^{y^{1/\beta}}
\,y^{\frac{\alpha+\gamma+1-\mu}{\beta}}}
{(y^{\alpha/\beta}-x)(y^{\gamma/\beta}-z)}\\
&+
\frac{1}{2\pi i}\,
\frac{1}{\alpha}
\int_{\omega(\varepsilon_\alpha;\eta_\alpha)}
\frac{
e^{t^{1/\alpha}}
\,t^{\frac{\beta+\gamma+1-\mu}{\alpha}}
}
{
(t-x)(t^{\beta/\alpha}-y)(t^{\gamma/\alpha}-z)
}
\,dt.\end{aligned}\end{equation}
Under the transformation
$
s=t^{\beta\gamma},
$
the Hankel contour
$
\omega(\varepsilon;\eta)
$
is mapped onto the Hankel contour
$
\omega(\varepsilon_\alpha;\eta_\alpha).
$ Choose an arbitrary number $\widetilde{\varepsilon}_{\alpha}$ satisfying
$\widetilde{\varepsilon}_{\alpha}>|x|$. Then
$
x\in\Omega^{(-)}(\widetilde{\varepsilon}_{\alpha};\eta_\alpha).
$
Furthermore, setting
$
\widetilde{\varepsilon}_{\beta}
=\widetilde{\varepsilon}_{\alpha}^{\frac{\beta}{\alpha}},
\,\,
\widetilde{\varepsilon}_{\gamma}
=\widetilde{\varepsilon}_{\alpha}^{\frac{\gamma}{\alpha}},
$
it follows that
$
y\in\Omega^{(-)}(\widetilde{\varepsilon}_{\beta};\eta_\beta),
\,\,
z\in\Omega^{(-)}(\widetilde{\varepsilon}_{\gamma};\eta_\gamma).
$ Hence, applying the integral representation \eqref{3.17}, we arrive at
\begin{equation}\label{3.18}\begin{aligned}
E_{\alpha,\beta,\gamma}(x,y,z;\mu)
&=\frac{1}{\gamma}\frac{
e^{z^{1/\gamma}}
\,z^{\frac{\alpha+\beta+1-\mu}{\gamma}}}
{(z^{\alpha/\gamma}-x)(z^{\beta/\gamma}-y)}+\frac{1}{\beta}\frac{
e^{y^{1/\beta}}
\,y^{\frac{\alpha+\gamma+1-\mu}{\beta}}}
{(y^{\alpha/\beta}-x)(y^{\gamma/\beta}-z)}\\
&+
\frac{1}{2\pi i}\,
\frac{1}{\alpha}
\int_{\omega(\widetilde{\varepsilon}_{\alpha};\eta_\alpha)}
\frac{
e^{t^{1/\alpha}}
\,t^{\frac{\beta+\gamma+1-\mu}{\alpha}}
}
{
(t-x)(t^{\beta/\alpha}-y)(t^{\gamma/\alpha}-z)
}
\,dt.\end{aligned}\end{equation}
Moreover, if
$
\varepsilon_\alpha<|x|<\widetilde{\varepsilon}_{\alpha},
$
then
$
|\arg x|<\eta_\alpha,
$
and, by Cauchy's theorem,
\begin{equation}\label{3.19}\begin{aligned}
E_{\alpha,\beta,\gamma}(x,y,z;\mu)
&=
\frac{1}{2\pi i}\,
\frac{1}{\alpha}
\int_{\omega(\widetilde{\varepsilon}_{\alpha};\eta_\alpha)-\omega(\varepsilon_{\alpha};\eta_\alpha)}
\frac{
e^{t^{1/\alpha}}
\,t^{\frac{\beta+\gamma+1-\mu}{\alpha}}
}
{
(t-x)(t^{\beta/\alpha}-y)(t^{\gamma/\alpha}-z)
}
\,dt\\
&=\frac{1}{\alpha}\frac{
e^{x^{1/\alpha}}
\,x^{\frac{\beta+\gamma+1-\mu}{\alpha}}}
{(x^{\beta/\alpha}-y)(x^{\gamma/\alpha}-z)}\end{aligned}
\end{equation}
Combining \eqref{3.19}, \eqref{3.18}, \eqref{3.17}, and \eqref{3.16}, we obtain
 the following integral representation 
\begin{equation}\label{3.20}\begin{aligned}
&E_{\alpha,\beta,\gamma}(x,y,z;\mu)\\
&=\frac{1}{\alpha}\frac{
e^{x^{1/\alpha}}
\,x^{\frac{\beta+\gamma+1-\mu}{\alpha}}}
{(x^{\beta/\alpha}-y)(x^{\gamma/\alpha}-z)}+\frac{1}{\beta}\frac{
e^{y^{1/\beta}}
\,y^{\frac{\alpha+\gamma+1-\mu}{\beta}}}
{(y^{\alpha/\beta}-x)(y^{\gamma/\beta}-z)}+\frac{1}{\gamma}\frac{
e^{z^{1/\gamma}}
\,z^{\frac{\alpha+\beta+1-\mu}{\gamma}}}
{(z^{\alpha/\gamma}-x)(z^{\beta/\gamma}-y)}\\
&+
\frac{1}{2\pi i}\,
\frac{1}{\alpha\beta\gamma}
\int_{\omega(\varepsilon;\eta)}
\frac{
e^{s^{1/(\alpha\beta\gamma)}}
\,s^{\frac{\alpha+\beta+\gamma+1-\mu}{\alpha\beta\gamma}-1}
}
{
(s^{1/\beta\gamma}-x)(s^{1/\alpha\gamma}-y)(s^{1/\alpha\beta}-z)
}
\,ds.\end{aligned}\end{equation}
Thus Lemma \ref{l3.2} is proved. \qed
\end{proof}

\end{document}
